\chapter{数据集、场景复杂性与评测协议}
\label{cha:data_protocol}

\section{数据集与任务定义}

\subsection{任务与数据组织}

本文实验对象是电梯轿厢内的 \textit{person + ebike} 双目标检测。数据来自实际轿厢监控视频，抽帧后由人工使用 LabelImg 按 YOLO 格式标注，类别为 \textit{person} 和 \textit{ebike}。训练集约 3600 张标注图像，验证集固定为 720 张图像；遇到遮挡目标或局部可见目标时，标注采用可见区域外接矩形。图像验证集负责给出检测指标，代表性视频片段用于观察连续画面中的输出变化，板端运行记录则用于拆分阈值行为、运行状态和分段耗时。

实验设计围绕第一章的问题链展开，但不把指标堆成单纯列表。需要回答的核心问题包括：模型进入板端后能否跑完验证集，\textit{person} 与 \textit{ebike} 的分类别结果是否符合任务语义，置信度阈值、NMS 和应用层清理会不会改变安全输出，代表性视频片段中检测框是否连续，以及 batch 图像验证耗时、单帧核心处理耗时和模型执行阶段耗时能否区分清楚。第五章的结果组织按这些问题展开。

\subsection{场景复杂性因素}

电梯场景的困难之处来自封闭空间、固定视角和安全响应要求的叠加。轿厢内常见的高亮区域、门框结构、局部目标、近距离重叠、尺度变化和暗部区域都会改变目标轮廓、类别分数或候选框排序。这些因素在本文中作为误差解释维度，服务于阈值、NMS、后处理和视频稳定性实验中的现象分析。

本文采用的证据来自几类材料：720 张图像验证集、既有板端检测输出、NMS 参数实验、后处理消融、batch 图像复查、连续视频分段计时和代表性视频片段。图像验证集给出总体精度、分类别精度和 batch 验证耗时；阈值后验重算观察同一批输出在不同安全阈值下的变化；NMS 与后处理实验用来区分模型输出和部署规则对候选框的影响；分段计时拆开 batch 端到端耗时、单帧核心处理和模型执行阶段；视频片段用于观察连续帧中的漏检、闪烁和边框抖动。上述证据共同回答一个问题：同一 YOLO 模型迁移到 OM 并在板端运行后，检测行为是否仍能被解释和复查。

场景因素在本文中保持为误差解释维度。它们与图像级检测指标、板端运行记录一起使用，既避免论文主线从双目标检测发散开，也保留了电梯场景相对于通用检测任务的特殊性。

\section{评测协议设计}

\subsection{评测指标}

本文按用途组织指标，而不是简单分成若干并列类别。基础识别能力由 precision、recall、F1 和 mAP@0.5 描述；\textit{person} 与 \textit{ebike} 的 TP、FP、FN、recall 和 mAP@0.5 用来检查分类别语义；多阈值 F1、类别召回、NMS 参数下的 FP/FN 变化和后处理消融，用来观察部署参数对最终输出的影响。运行时部分则记录成功处理样本数、batch 图像验证链路端到端单图处理耗时和分段计时。跨阶段 IoU 漂移保留为更细粒度的方法设计，不作为本章主要实测结果。

\subsection{实验协议}

为保证结果可比较，图像级指标全部在同一验证集上统计，并统一保留三位小数。阈值实验只改变置信度阈值，其余参数保持不变；NMS 实验只改变 IoU 阈值；后处理消融只改变电动车清理策略；视频片段实验只改变平滑窗口；分段计时固定模型、二进制和阈值设置，分别记录 batch 端到端、单帧核心处理和模型执行阶段。本文并不追求实验数量堆叠，而是让每个实验只回答一个主要问题。

\subsection{场景复杂性与部署误差的关系}

场景复杂性会影响模型迁移后的结果解释。高亮区域和暗部区域常先影响类别置信度，在阈值扫描中表现为召回率波动；轿厢边界和门框结构更容易干扰边框回归与 NMS，进而拉低 IoU；近距离重叠则可能同时改变类别排序和候选框抑制关系，使不同阶段的输出数量不一致。

因此，实验分析需要把指标变化放回场景中理解。mAP@0.5 下降时，先判断问题来自类别分数还是边框定位；视频片段出现闪烁时，再结合阈值结果判断是否存在置信度边界问题。


\section{数据可追溯性与实验设计}

\subsection{数据集划分与可追溯性}

实验结果要能复查，数据划分必须先说明清楚。本文图像验证集固定为 720 张，用于报告板端总体指标和分类别指标。论文正文只交代验证集规模、标注类别和使用目的，不展开本地文件组织细节，以免把工程路径写成论文证据。

可追溯性主要体现在样本数和统计口径上。实验表格中的样本数需要与指标来源一致；\textit{person} 与 \textit{ebike} 的真实目标数、预测数、TP、FP、FN 要能对应到评测结果；成功样本数、失败样本数和端到端耗时则用于说明板端验证是否按预期执行。

视频片段采用检测保持率、闪烁窗口、连续无输出片段和分段计时等机器侧统计描述输出波动。这样可以让指标和实验材料相匹配。


\subsection{跨阶段一致性实验设计}

跨阶段一致性实验的理想形式是选取同一批输入样本，分别在训练阶段、ONNX 导出阶段、OM 编译阶段和板端部署阶段运行，再对输出结果进行逐样本匹配。匹配时先按照类别和 IoU 阈值确定同一目标，然后统计类别一致率、平均边框 IoU、置信度差异和最终检测数量差异。该实验能够直接衡量模型迁移过程是否改变检测行为。

在实际执行时，应避免同时改变多个变量。输入图像、预处理方式、置信度阈值和 NMS 阈值需要固定，唯一变化的是推理阶段。若发现输出差异，再根据第三章的误差定位表检查类别契约、坐标契约或后处理契约。这样，跨阶段实验不仅给出结果，还能帮助定位问题来源。

本文已经积累了板端全量验证、参数敏感性、后处理消融、视频连续输出和分段计时证据。训练阶段、ONNX 导出阶段、OM 编译阶段和板端部署阶段的逐样本统一对比可在此基础上进一步展开，本节给出其分析方式，用于衔接后续更细粒度的迁移分析。


\subsection{实验结果呈现原则}

为了保证实验结果可信，本文区分实测结果、后验重算结果和方法性评价设计。已有完整证据支撑的内容，如 720 张图像的板端验证结果、NMS 参数敏感性、后处理消融、代表性视频稳定性和分段计时结果，可以写入实验结果并进行结论分析；基于既有检测输出重新筛选得到的阈值敏感性结果，作为补充实验处理；跨阶段逐样本对齐则保留为进一步细粒度分析的评价方法。

上述做法与实验研究对证据分层的基本要求一致。工程实践中存在日志、脚本、视频和人工观察等多种材料，只有区分证据强度，才能避免把经验观察混同为定量结论。本文用图像验证集计算检测指标，用参数敏感性和后处理消融解释部署规则影响，用视频片段辅助观察稳定性，再用场景复杂性描述解释误差来源。

经过这样的划分，实验结果、后验重算和方法设计各有作用，部署一致性分析可以在实测证据基础上逐层展开。

\section{本章小结}

本章给出了本文的数据组织方式、场景复杂性因素、评测指标体系和实验协议。场景复杂性描述用于误差解释，分析重点放在 YOLO 双目标检测模型迁移到 HongOU PI（Hi3516DV500）开发板后的部署行为是否对齐。
